\section{A Influência da Iluminação na Generalização dos Modelos}

Um dos achados mais significativos deste estudo foi a variação de performance entre as versões V1 e V10, motivada primariamente pela alteração nas condições físicas de iluminação durante a captura. A superioridade da V10 reforça o aforismo \textit{garbage in, garbage out} (lixo entra, lixo sai) no contexto industrial.

Observou-se que a padronização e o espectro da fonte de luz têm um impacto mais direto na acurácia do que o ajuste fino de múltiplos hiperparâmetros. Modelos treinados sob iluminação otimizada demonstraram menor entropia na função de perda e convergência mais estável. Para o setor têxtil, este resultado sugere que o investimento em infraestrutura física de captura (câmaras de luz controladas) oferece um retorno sobre investimento (ROI) em precisão algorítmica superior ao desenvolvimento de modelos excessivamente complexos que tentam compensar capturas ruidosas.
